\cotefile{156-9}
\scanpage{74}
\mspage{70}
Ceci posé, \deleted{\unsure{les} $\beta_{\unsure{\ast|\times}}$} \unsure{on aura|on a}
\[
\beta_{\ast}(x')=C\bigl(\underbrace{\bar{\beta}(x')}_{\substack{\uparrow\\ \mathbb{\Sigma}_{E}}}\bigr)=\operatorname{cosupp}_{E}^{\circ}\bigl(\bar{\beta}(x')\bigr)
\]
et l'application $\beta$ \unsure{de|en} \unsure{(ii)|(vi)} satisfait la condition \unsure{(xx)|(ix)} \unsure{ssi} $\beta_{\ast}$ est \unsure{associée|associé} à $\beta$, ou encore
\[
\operatorname{supp}_{E}^{\circ}\beta(x')=C\bigl(\bar{\beta}(x')\bigr)
\]
ou aussi
\[
\operatorname{cosupp}^{\circ}\unsure{C|(}\beta(x')=\bar{\beta}(x').
\]

\uline{Définition} Les $\alpha_{\ast}$ satisfaisant les conditions équivalentes \unsure{(i)} \unsure{(bis)|(ii)} \unsure{(bis$'$)|(ii$'$)} sont appelées correspondances disjonctives \deleted{\unsure{entre}} $\underline{E},\underline{E}'$. La correspondance $\beta_{\ast}$ \unsure{de|du} (i) \inserted{(qui est aussi une corr. disj.)} \unsure{est} appelée \deleted{\unsure{la}} corr. disjonctive \deleted{\uline{\unsure{associée}}} \uline{transposée} de $\alpha_{\ast}$.

\uline{Scholie} On trouve \deleted{ainsi} une bijection
\[
\alpha_{\ast}\mathrel{\unsure{\vdash|\mapsto}}{}^{t}\alpha_{\ast}:\operatorname{Corrdis}(\underline{E},\underline{E}')\longrightarrow\operatorname{Corrdis}(\underline{E}',\underline{E})
\]
et on a
\[
\left\{\begin{aligned}
{}^{t}({}^{t}\alpha_{\ast})&=\alpha_{\ast}\\
{}^{t}({}^{t}\beta_{\ast})&=\beta_{\ast}
\end{aligned}\right.
\]
% archive sheet number: 73
\scanpage{75}
\mspage{71}
De plus, l'on a \deleted{alors}, pour $A\subset E$
\[
\underset{\substack{\vrule height .45ex width .08ex\rule{7.2em}{.08ex}\vrule height .45ex width .08ex\\
\Vert\,\mathrm{def}\\
\operatorname{Supp}_{E'}^{\circ}\bigl(\bar{\alpha}(\operatorname{supp}_{E}^{\circ}(A)\bigr)}}{\alpha_{\unsure{\ast|\times}}\bigl(\operatorname{supp}_{E}^{\circ}(A)\bigr)}
=\deleted{\unsure{s}}\operatorname{Supp}_{E'}^{\circ}\bigl(\bar{\alpha}(\deleted{\unsure{E|A}}\unsure{A|E})\bigr)
\tag*{\makebox[0pt][r]{\makebox[\displaywidth][l]{\unsure{(xxx)|($\times\times\times$)}}}}
\]

\uline{NB} J'ignore si la validité de cette formule, pour une \deleted{$\alpha_{\ast}$} quasi-corr. disjonctive donnée $\alpha_{\ast}$, indique que c'est une corr. disjonctive

En analogie avec le cas des disjoncteurs triviaux, où on distingue le cas \unsure{d'}une correspondance qui est une application, ou \unsure{transposée|transposé} d'une application, je dirai que $\alpha_{\ast}$ est un \uline{morphisme} si $ {}^{t}\alpha_{\ast}=\beta_{\ast}$ \unsure{a} \deleted{\unsure{la}} les propriétés suivantes : il commute à $C$ (donc aux $\operatorname{Inf}$ quelconques)

% archive sheet number: 74
\scanpage{76}
\mspage{72}
\uline{10 juillet}

Je reprends \unsure{du début|dès début} le formalisme des disjoncteurs, spatios.

1. \uline{\unsure{Disjoncteurs|Disjoncteur}}, \inserted{\uline{\unsure{supports|support}}.} : Ensemble $E$ muni d'une relation $\mathrel{|\unsure{\circ|0}|}$ \inserted{(dite \uline{rel. de disjonction})}, on \deleted{symétrique. \unsure{O}} \inserted{suppose}

\[
\begin{aligned}
&\text{\uline{Disj.} 1}\qquad x\mathrel{|\circ|}y\ \unsure{\longrightarrow|\text{est}}\ \text{symétrique}\\
&\text{\uline{Disj.} 2}\qquad x\mathrel{|\circ|}x\Longrightarrow x\mathrel{|\circ|}y\quad\forall y\in E
\end{aligned}
\]

% slanted
\inserted{\uline{Structure de disjonction}.} \inserted{On désigne par $E_{0}$ (``\deleted{\unsure{singulariseur|singulariser}}\unsure{singularité|défectuosité}'' du disj.) l'ens. \inserted{des $x$}}
Disjoncteur antiréflexif : quand $\mathrel{|\circ|}$ est antiréflexif (ce qui implique Disj. 2), i.e. $E_{0}=\emptyset$.

\uline{Relation à notion de graphe} (strict). On associe à une str. de disj. $\mathrel{|\circ|}$ sur $E$, une structure de graphe, en prenant comme ens. d'arêtes
\[
A\subset\unsure{\mathcal{P}|\mathrm{P}}_{2}(E)
\]
l'ens. suivant
\[
A=\bigl\{\{x,y\}\in\mathcal{P}_{2}(E)\mid x\mathrel{\overline{|\circ|}}y\bigr\}.
\]

On récupère \inserted{presque} la structure de disjoncteur à partir de la structure de graphe, \deleted{\unsure{en partant}} \inserted{il faut de} plus se donner la partie $E_{0}$ de l'ens.
\[
\deleted{x\mathrel{|\circ|}y\Longleftrightarrow\bigl(\{x,y\}\in A,\ \text{ou }x=y\text{ et }x\ \text{\uline{sommet}}}
\]
des \deleted{\unsure{points}} \inserted{sommets} isolés du graphe, \deleted{\unsure{pour} \unsure{élément} $E_{0}$ \unsure{isolé du graphe} $(E,A)$\unsure{)}}.

\uline{La structure de disjoncteur antiréflexif}$^{\ast}$ \uline{équivaut à celle de graphe strict}
% slanted
\margin{$^{\ast}$ C'est le cas le plus important, bien sûr.}

(On trouve ainsi une corr. bi-univ. entre structures de disjonction sur $E$, et structures de graphe strict, i.e. parties de $\mathcal{P}_{2}(E)$.

La structure de disjoncteur donne un \unsure{formalisme} de \inserted{$\operatorname{cosupp}^{\circ}$ ou} $\operatorname{cosupp}_{E}^{\circ}$ (ou $\operatorname{cosupp}_{\underline{E}}^{\circ}$, \unsure{où} $\deleted{\unsure{E}}\underline{E}=(E,\mathrel{|\circ|}_{E})$) et $\operatorname{supp}_{E}^{\circ}=\operatorname{cosupp}_{E}^{\circ}\circ\operatorname{cosupp}_{E}^{\circ}$ \deleted{\unsure{et}} de $\mathcal{P}(E)$, \unsure{d'est|dont|d'où|c'est}
% archive sheet number: 75
\scanpage{77}
\mspage{73}
partie $\mathbb{\Sigma}_{\underline{E}}$ (ou $\mathbb{\Sigma}_{E}$) de $\unsure{\mathcal{P}|\mathrm{P}}_{E}$, stable par $\bigcap_{i}$ \unsure{quelconque|quelconques}, ayant comme \unsure{pl. grand|pl. grand.} \unsure{élémt|élément} $E$, comme plus petit élément le \unsure{sous-ensemble} $E_{0}$. \deleted{\unsure{Il}}\unsure{Elle|Il} a des $\operatorname{Inf}$ quelc. (les $\bigcap_{i}$) et $\operatorname{Sup}$ quelc. (en général $\deleted{\unsure{\not\leq|\neq}}\deleted{\unsure{les}}\ne\bigcup_{i}$), et une anti-involution
\[
C\deleted{\unsure{A}}\unsure{S|A}\overset{\mathrm{def}}{=}\operatorname{cosupp}_{\deleted{\unsure{E}}}^{\circ}A
\qquad (A\in\mathbb{\Sigma}_{E}\subset\mathcal{P}(E)).
\]

Application canonique
\[
x\longmapsto\sigma(x)=\operatorname{supp}_{E}^{\circ}(\{x\}):E\xrightarrow{\sigma}\mathbb{\Sigma}_{E}
\]
% margin note written sideways
\margin{On a $x\mathrel{|\unsure{\circ|0}|}y\Longleftrightarrow\sigma x\mathrel{|\circ|}\sigma y$ où pour $S,T\in\mathbb{\Sigma}_{\unsure{E}}$, plus \unsure{génér.} pour $A,B\in\mathcal{P}(E)$, on pose $A\mathrel{|\circ|}B\Longleftrightarrow x\mathrel{|\circ|}y\quad\forall x\in A,y\in B$. \unsure{Or.|On.} \deleted{Si $S,T\in\mathbb{\Sigma}_{E}$, \unsure{on}} $A\mathrel{|\circ|}B\Longleftrightarrow A\subset\operatorname{cosupp}^{\circ}B\Longleftrightarrow B\subset\operatorname{cosupp}^{\circ}A\Longleftrightarrow\operatorname{supp}^{\circ}(A)\subset\operatorname{cosupp}^{\circ}B\Longleftrightarrow\operatorname{supp}^{\circ}B\subset\operatorname{cosupp}^{\circ}A$. Dans $\mathbb{\Sigma}_{E}$, la relation $\mathrel{|\circ|}$ \deleted{\unsure{de}} résulte donc de la structure $\leq$ et $C$.}

\uline{Disjoncteur} \uline{trivial} : quand
\[
x\mathrel{|\circ|}y\Longleftrightarrow x\ne y
\]
donc cela correspond \unsure{:|à} $A=\emptyset$, structure de graphe \unsure{discrète|discrèt}. On a alors
\[
\mathbb{\Sigma}_{E}=\mathcal{P}(E).
\]

\uline{NB} Pour toute partie $A$ de $E$, on a
\[
\begin{aligned}
\operatorname{cosupp}_{E}^{\circ}(A)&=\operatorname{cosupp}_{E}^{\circ}(A\cup E_{0})\supset E_{0},\qquad\text{d'où}\\
\deleted{\unsure{\mathrm{c}}}\operatorname{Supp}_{E}^{\circ}(\unsure{M|A})&=\operatorname{supp}_{E}^{\circ}(A\cup E_{0})\unsure{>|\supset} E_{0},
\end{aligned}
\]
et par suite aussi
\[
\begin{aligned}
\operatorname{cosupp}_{E}^{\circ}(A)&=\operatorname{cosupp}^{\circ}(A^{\unsure{\ast|\prime}})\\
\operatorname{supp}_{E}^{\circ}(M)&=\operatorname{supp}^{\circ}(A^{\ast})
\end{aligned}
\]
où
\[
A^{\ast}=A\setminus E_{0}\cap A.
\]
On en conclut que l'on a un iso. (respectant structure d'ordre et $C$)
\[
\begin{aligned}
\mathbb{\Sigma}_{E}\mathrel{\overset{\sim}{\longleftrightarrow}}\mathbb{\Sigma}_{E^{\ast}}\qquad &\deleted{\unsure{A}}S\longmapsto S\setminus E_{0}\\
&S^{\ast}\mathbin{\unsure{\sqcup|\cup}}E_{0}\mathrel{\longleftarrow\mkern-9mu\vert} S^{\ast}.
\end{aligned}
\]
% archive sheet number: 76
\scanpage{78}
\mspage{74}
Cela montre que pour la construction

2. \uline{Spatios} Deux façons de \unsure{le|les} décrire.

\deleted{\unsure{Pr|P}}\uline{Définition 1} Un spatios est un ensemble $\Sigma$, muni \unsure{de$'$|de} \deleted{deux} relation\deleted{s} $\leq$, $\deleted{\unsure{\mathrel{|\circ|}|\mathrel{\#}}}$ et d'un endomorphisme $C$ (relation aussi, si on veut !), satisfaisant les conditions

\uline{Spat 1} La relation $\leq$ est une \uline{relation d'ordre}, admettant des $\operatorname{Sup}$ et des $\operatorname{Inf}$ quelconques (donc un plus petit élément, noté \inserted{$0$ ou} $0_{\Sigma}$, et un plus grand élément, noté \inserted{$1$ ou} $1_{\Sigma}$).

\uline{Spat 2} $C$ est une \uline{involution} ($C\deleted{\unsure{x}}(Cx)=x$) et \deleted{et} ``renverse la relation d'ordre'' :
\[
\deleted{\unsure{C}}x\leq y\Longleftrightarrow Cy\leq Cx.
\]

\deleted{\uline{\unsure{Spat 3}} $\forall x$}
\[
\left[\begin{aligned}
&\text{\uline{Cor} 1}\quad C(\operatorname{Sup}_{i}x_{i})=\operatorname{Inf}_{i}Cx_{i},\quad C(\operatorname{Inf}_{i}x_{i})=\operatorname{Sup}_{i}Cx_{i}.\\
&\text{\uline{Cor. 2}}\quad C0=1,\quad C1=0
\end{aligned}\right.
\]

\deleted{A un Spatios est \unsure{associé|associée} une structure de disjonction canonique, en \unsure{posant}}
\[
\deleted{x\mathrel{|\unsure{\circ|0}|}y\mathrel{\overset{\mathrm{def}}{\Longleftrightarrow}}x\leq Cy\quad(\Longleftrightarrow y\leq Cx\ \text{par Spat 2}).}
\]

\uline{Spat 3} $\forall x\in\Sigma$, on a
\[
x\wedge Cx=0_{\deleted{\unsure{E}}\Sigma}
\]
\unsure{ce|ou} qui \unsure{équiv|équiv.} par Spat 2 (Cor. 1 et 2) :
\[
x\unsure{:}\vee Cx=1_{\Sigma}.
\]
% archive sheet number: 77
\scanpage{79}
\mspage{75}
\deleted{\unsure{À|A}} Un spatios est \deleted{associé} \inserted{muni} d'une structure de disjonction canonique :
\[
x\mathrel{|\unsure{\circ|0}|}y\mathrel{\overset{\mathrm{def}}{\Longleftrightarrow}}x\leq Cy\quad(\Longleftrightarrow y\leq Cx,\ \text{par Spat 2}).
\]
\unsure{Notons|Noter|Observons} que
\[
x\mathrel{|\circ|}y,\ x'\leq x,\ y'\leq y\Longrightarrow x'\mathrel{|\circ|}y'.
\]
% slanted
\inserted{\uline{NB} Si $x=\operatorname{Sup}x_{i}$, $y=\operatorname{Sup}y_{j}$, $x_{i}\mathrel{|\circ|}y_{j}\quad\forall i,j$, alors $x\mathrel{|\circ|}y$ (et inversement)}

\deleted{La structure de disjonction permet de récupérer la structure $\leq,C$, grâce à la}

\deleted{\uline{Proposition} Soit $\Sigma$ un spatios, considérons}

\unsure{1|l'} \uline{Exemple de Spatios} Si $\underline{E}=(E,\mathrel{|\circ|})$ est un disjoncteur, $(\mathbb{\Sigma}_{\underline{E}},\subset,C)$ est un spatios. Les \deleted{\unsure{relations}} propriétés Spat 1, Spat 2 ont été dégagées, elles ne dépendent que du fait que $\mathrel{|\circ|}_{E}$ soit symétrique \inserted{(Disj 1)}, sans plus. La propriété Spat 3 résulte de (Disj. 2), qui \unsure{implique|signifie}, pour $\forall A\in\unsure{\mathcal{P}|\mathrm{P}}(E)$
\[
A\cap\operatorname{cosupp}_{E}^{\circ}(A)\subset E_{0}
\]
donc, si $A\supset E_{\unsure{0}}$, et en particulier si $A\unsure{\in|\subset}\mathbb{\Sigma}_{E}$
\[
A\cap\operatorname{cosupp}_{E}^{\circ}(A)=E_{0}.
\]

% slanted
\margin{\uline{Ensemble ordonné disjonctif} ou \uline{disjoncteur ordonné}. $(E,\overset{\circ}{\ll},\mathrel{|\circ|})$ avec $x\mathrel{|\circ|}y,\ x'\mathrel{\overset{\circ}{\ll}}x,\ y'\mathrel{\overset{\circ}{\ll}}y\Longrightarrow x'\mathrel{|\circ|}y'$. Ex : Tout spatios est un disjoncteur ordonné. Si $(E,\mathrel{\overset{\circ}{\ll}},\mathrel{|\circ|})$ est un disjoncteur ordonné, $\sigma_{E}:E\longrightarrow\mathbb{\Sigma}_{\underline{E}}$ est \uline{croissante}.}

% slanted
\margin{\uline{NB} Tout disj. \deleted{\unsure{ordonné}} est muni d'une relation d'ordre \uline{canonique}, qui en fait un disj. ordonné (la \unsure{moins} fine de toutes celles qui sont \unsure{comp.} avec la str. de disjonction) : \unsure{à savoir|c'est|à dire} l'image inverse de la structure d'ordre de $\mathbb{\Sigma}_{E}$.}

Ainsi, la structure de disjonction sur un \unsure{disjoncteur|ensemble} quelconque, \unsure{se|on} déduit de celle sur le \deleted{disjoncteur} spatios associé $\mathbb{\Sigma}_{\underline{E}}$, via l'application canonique
\[
\sigma_{\underline{E}}:E\longrightarrow\mathbb{\Sigma}_{\underline{E}}.
\]
% archive sheet number: 78
\scanpage{80}\mspage{76}
La structure $\mathrel{|\unsure{\circ|0}|}$ sur un Spatios permet de récupérer la structure $\leq,C$ du spatios, \unsure{grace|grâce} à la proposition suivante :

% slanted
\margin{Décrire la relation d'ordre \unsure{can.} d'un disjoncteur : $x\leq y\Longleftrightarrow\forall z,\ z\mathrel{|\circ|}y\Longrightarrow z\mathrel{|\circ|}x$}

\uline{Proposition 1} Soit $(\deleted{\unsure{S}}\unsure{\mathbb{\Sigma}|\Sigma},\leq,C)$ un spatios, $\mathrel{|\circ|}$ la structure de disjonction associée. Alors l'application canonique
\[
\sigma:\Sigma\longrightarrow \mathbb{\Sigma}_{\Sigma}
\]
\uline{est} \deleted{\unsure{bijec|bij}} \uline{un isomorphisme de spatios}.

L'\deleted{\unsure{appl}}isomorphisme inverse est donné par
\[
\sigma^{-1}:S\longmapsto\operatorname{Sup}S.
\]

% slanted
\margin{\uline{NB} La relation d'ordre est exprimée en termes de $\mathrel{|\circ|}$ comme ci-dessus. D'autre part, $Cx=\underset{\substack{x'\in\Sigma\\\text{t.q. }x'\mathrel{|\circ|}x}}{\operatorname{Sup}}x'$}

\uline{Scholie} Pour un ensemble donné $\Sigma$, il y a corr. 1-1 \unsure{dentre|entre} les deux ensembles suivants :

(i) L'\unsure{ens.|Ens.} \inserted{$\mathscr{S}$} des structures de spatios \inserted{$(\leq,C)$} sur $\deleted{\unsure{E}}\Sigma$.

(ii) L'\unsure{ens|en} \inserted{$\mathscr{D}$} des structures de disjonction \inserted{$\mathrel{|\circ|}$} sur $\Sigma$, telles que l'application canonique
\[
\Sigma\longrightarrow\mathbb{\Sigma}_{\underline{\Sigma}}\qquad
(\text{où }\underline{\Sigma}\unsure{=|\simeq}(\Sigma,\mathrel{|\circ|}))
\tag*{\makebox[0pt][r]{\makebox[\displaywidth][l]{($\unsure{\ast|\times}$)}}}
\]
soit bijective.

\deleted{\unsure{A l'appli|A l'appl}} La bijection $\mathscr{S}\overset{\sim}{\longrightarrow}\mathscr{D}$ \deleted{\unsure{est donnée}} \inserted{\unsure{associe|associé}, à} \deleted{par la} une structure de spatios, la structure de disj. associée. La bijection inverse associe à une structure de disjonction $\in\mathscr{D}$, la structure de spatios sur $\Sigma$ \unsure{déduit|déduite} par tr. de str. de la bijection ($\ast$).

% archive sheet number: 79
\scanpage{81}\mspage{77}
Ainsi, les spatios peuvent être considérés comme des disjoncteurs particuliers, qu'on pourrait appeler \uline{disjoncteurs} \uline{canoniques}. \deleted{\unsure{Si}}

Si $(E,\inserted{\mathrel{|\unsure{\circ|0}|}})$ est un disjoncteur quelconque, on aimerait considérer $ \mathbb{\Sigma}_{E}$, muni de
\[
\sigma_{E}:E\unsure{-|\longrightarrow}\mathbb{\Sigma}_{E},
\]
comme un ``disjoncteur canonique'' (ou \unsure{Spatios|spatios}) \uline{enveloppant} de $E$. \unsure{Je voudrais|Je voudraits} déjà caractériser, à is. unique près, le spatios $\mathbb{\Sigma}_{E}$ et l'application $\sigma$.

\uline{Proposition 2} \unsure{Soient|Soit} $(E,\mathrel{|\circ|})$ un disjoncteur, $(\deleted{\unsure{S|\mathscr{S}}}\Sigma,\leq,C)$ un spatios, et
\[
\alpha:E\longrightarrow\Sigma
\]
une application. Pour qu'il existe un diagramme commutatif
\diagram{\begin{tikzcd}[ampersand replacement=\&,row sep=0.6em,column sep=3em]
{} \& \deleted{\unsure{E|S}}\Sigma \\
\unsure{,}\,E \arrow[ur,"{\deleted{\unsure{\alpha|\varphi}}\sigma}"] \arrow[dr,"{\sigma_E}"] \& {} \\
{} \& \mathbb{\Sigma}_E \arrow[uu,leftrightarrow,"{\rotatebox{90}{$\sim$}\ \substack{\deleted{\unsure{\alpha|\varphi}}\\ \unsure{\varphi|1}}}"{swap}] \& {,}
\end{tikzcd}}
où $\varphi$ est un iso, il f. et s. que $\sigma$ satisfasse les \inserted{deux} conditions suivantes
\[
\forall x,y\in E,\qquad x\mathrel{|\circ|}y\Longleftrightarrow
\deleted{\unsure{\alpha|\varphi}}\sigma x\mathrel{|\circ|}\deleted{\unsure{\alpha|\varphi}}\sigma y
\tag*{\makebox[0pt][r]{\makebox[\displaywidth][l]{a)}}}
\]
% archive sheet number: 80
